\section{Method}
\label{sec:method}

\model{} predicts atomic coordinates and a unit cell from a molecular specification, a formula-unit count $Z$, and optional auxiliary information using a conditional coordinate--lattice flow. We describe the selected \model{} configuration below, with design choices justified by the controlled studies in \Cref{sec:recipe}.

\subsection{Problem Formulation}
\label{sec:method-problem}

\textbf{Crystal representation.}
We represent a crystal with $N$ atoms by atom types
$\mathbf{A}=(A_1,\ldots,A_N)\in\mathcal{A}^{N}$, Cartesian coordinates
$\mathbf{X}=(\mathbf{x}_1,\ldots,\mathbf{x}_N)^{\top}\in\mathbb{R}^{N\times3}$, and a
lattice cell $\mathbf{L}\in\mathbb{R}^{3\times3}$ whose rows are lattice
vectors. Integer translations along these lattice vectors define the corresponding
infinite periodic crystal.

Since the same lattice can be represented by different basis vectors and global
rotations, we canonicalize the cell in two steps. First, we Niggli-reduce~\citep{krivy1976unified,grosse2004numerically} each
crystal to obtain a reduced basis $\widetilde{\mathbf{L}}$, reducing the ambiguity among equivalent lattice bases. Second, we remove the remaining global rotational
freedom by forming the rotation-invariant Gram matrix
$\widetilde{\mathbf{L}}\widetilde{\mathbf{L}}^{\top}$ and taking its unique
lower-triangular Cholesky factor with positive diagonal:
\begin{equation*}
\mathbf{L}=\operatorname{chol}\!\left(
\widetilde{\mathbf{L}}\widetilde{\mathbf{L}}^{\top}
\right)=
\begin{pmatrix}
e^{\ell_1} & 0 & 0 \\
\ell_2 & e^{\ell_3} & 0 \\
\ell_4 & \ell_5 & e^{\ell_6}
\end{pmatrix}.
\end{equation*}
We further parameterize the canonical cell by an unconstrained vector  $\boldsymbol{\ell}=(\ell_1,\ldots,\ell_6)\in\mathbb{R}^{6}$ following \citet{veljkovic2026crystalite}.

\textbf{Molecular crystal structure prediction (CSP).}
We formulate molecular CSP as conditional generation of Cartesian coordinates
$\mathbf{X}$ and the lattice latent $\boldsymbol{\ell}$ from a chemical
specification, a supplied formula-unit count $Z$, and optional auxiliary
information. The chemical specification consists of $K$ distinct molecular
graphs $\mathcal{G}=\{G_k\}_{k=1}^{K}$ and a stoichiometry vector
$\mathbf{r}=(r_1,\ldots,r_K)$, where $r_k$ denotes the number of copies of
component $G_k$ in one formula unit. For example, a binary $2{:}1$ cocrystal
has $\mathbf{r}=(2,1)$. Each molecular graph specifies atom types, bond types,
and formal charges. To be specific, if component $G_k$ contains $N_k$ atoms, a unit cell with $Z$ formula units
contains
\[
N = Z\sum_{k=1}^{K} r_k N_k
\]
atoms. Optional information $\mathcal{O}$ may include molecular templates,
stereochemical annotations, and a space-group label. A molecular template is a
reference 3D conformer, generated for example with RDKit~\citep{rdkit}, that
provides local molecular geometry without specifying crystal packing. Given the
complete condition $\mathcal{C}=(\mathcal{G},\mathbf{r},Z,\mathcal{O})$, the
generator models $p_{\theta}(\mathbf{X},\boldsymbol{\ell}\mid\mathcal{C})$.

\subsection{Variational Flow Matching for Molecular CSP}
\label{sec:method-vfm}

We use variational flow matching \citep[VFM;][]{eijkelboom2024variational}, which generalizes flow matching~\citep{lipman2023flow}
through a flexible choice of the
variational posterior~\citep{zaghen2025riemannian}.

\textbf{Training objective.} Let $(\mathbf{y}_1,\mathcal{C})$ denote a crystal state and its condition, and let
$\mathbf{y}_0\sim p_0$ be a prior sample.
We use the linear conditional path to interpolate between the prior and data endpoints:
\begin{equation}
\label{eq:method-flow-path}
\mathbf{y}_t=(1-t)\mathbf{y}_0+t\mathbf{y}_1.
\end{equation}

Given a noisy state $\mathbf{y}_t$,
VFM learns a variational posterior
$q_{\theta}(\mathbf{y}_1\mid\mathbf{y}_t,t,\mathcal{C})$
that approximates the true posterior by minimizing the expected negative log-likelihood,
\begin{equation}
\label{eq:method-vfm-general}
\mathcal{L}_{\mathrm{VFM}}(\theta)
=-\mathbb{E}_{\mathbf{y}_1,\mathbf{y}_t,t}\!\left[
\log q_{\theta}(\mathbf{y}_1\mid\mathbf{y}_t,t,\mathcal{C})
\right].
\end{equation}
The variational family determines the form of this objective.
We model the posterior as a fully factorized Laplace distribution with a fixed scale,
which yields a component-weighted $L^1$ loss on the predicted endpoints.

The model parameterizes the mean of the variational posterior,
which is equivalent to predicting the clean endpoint $\mathbf{y}_1$ from the noisy
state $\mathbf{y}_t$ and condition $\mathcal{C}$ as
$\widehat{\mathbf{X}}_1,\widehat{\boldsymbol{\ell}}_1
=\boldsymbol{\mu}_t^{\theta}\!\left(\mathbf{y}_t,\mathcal{C}\right)$.
The equation then simplifies to
\begin{equation}
\label{eq:method-vfm-loss}
\mathcal{L}_{\mathrm{VFM}}(\theta)
=\mathbb{E}_{\mathbf{y}_1,\mathbf{y}_t,t}\!\left[
\frac{\lambda_{\mathrm{coord}}}{3N}
\left\|\widehat{\mathbf{X}}_1-\mathbf{X}_1\right\|_1
+
\frac{\lambda_{\mathrm{lattice}}}{6}
\left\|\widehat{\boldsymbol{\ell}}_1-
\boldsymbol{\ell}_1\right\|_1
\right],
\end{equation}
where $\lambda_{\mathrm{coord}}$ and $\lambda_{\mathrm{lattice}}$ are loss weights
for the coordinate and lattice components, respectively.

\textbf{Sampling.}
After training $\boldsymbol{\mu}_t^\theta$, we generate samples by solving the ordinary differential equation (ODE) or the corresponding stochastic differential equation (SDE)~\citep{eijkelboom2024variational,albergo2025stochastic}:
\begin{equation}
d\mathbf{y}_t
=\mathbf{v}_t^\theta(\mathbf{y}_t)\,dt,
\qquad
d\mathbf{y}_t
=
\left[
\mathbf{v}_t^\theta(\mathbf{y}_t)
+\frac{g_t^2}{2}\mathbf{s}_t^\theta(\mathbf{y}_t)
\right]dt
+g_t\,d\mathbf{W}_t.
\end{equation}
For the linear path,
$\mathbf{v}_t^\theta(\mathbf{y}_t)
=(\boldsymbol{\mu}_t^\theta(\mathbf{y}_t,\mathcal C)-\mathbf{y}_t)/(1-t)$
and
$\mathbf{s}_t^\theta(\mathbf{y}_t)
=(t\,\mathbf{v}_t^\theta(\mathbf{y}_t)-\mathbf{y}_t)/(1-t)$.
Here, $g_t$ is the diffusion coefficient controlling the stochasticity, and
$\mathbf{W}_t$ is a standard Wiener process.

\subsection{Model Architecture}
\label{sec:method-architecture}

Our model $\boldsymbol{\mu}_t^{\theta}$ mainly employs a Pairmixer-based
architecture~\citep{ouyangzhang2025triangle} with trunk-head modularization~\citep{Abramson2024, bytedance2025protenix}. As shown in
\Cref{fig:method-architecture}, the model first constructs condition-only
single and pair representations, refines the pair representation with Pairmixer,
injects the noisy crystal state into the single representation, and processes it
with a Diffusion Transformer (DiT)~\citep{peebles2023scalable}. Separate output heads predict the coordinate and lattice endpoints. The complete architecture is specified
algorithmically in \Cref{app:model-architecture}.

% % Final crystal structure generator architecture.
% % Figure-local palette matches the architecture-ablation diagram.
% \definecolor{methodArchInputBg}{HTML}{E9F0F5}
% \definecolor{methodArchInputBlock}{HTML}{AFC9DA}
% \definecolor{methodArchInputLine}{HTML}{5D8198}
% \definecolor{methodArchTimeBg}{HTML}{FCF5E6}
% \definecolor{methodArchTimeBlock}{HTML}{F6DFB1}
% \definecolor{methodArchTimeLine}{HTML}{B7892D}
% \definecolor{methodArchPairBg}{HTML}{E3EEE9}
% \definecolor{methodArchPairBlock}{HTML}{C5DED1}
% \definecolor{methodArchPairLine}{HTML}{5F8D7A}
% \definecolor{methodArchTrunkBg}{HTML}{E7DCE3}
% \definecolor{methodArchHeadBg}{HTML}{F8EBE9}
% \definecolor{methodArchAttnBlock}{HTML}{EAC8B1}
% \definecolor{methodArchMlpBlock}{HTML}{CDB6C5}
% \definecolor{methodArchArrow}{HTML}{2D2D2D}
% \definecolor{methodArchFrame}{HTML}{30343B}

% \begin{figure}[t]
%   \centering
%   \resizebox{\textwidth}{!}{%
%   \begin{tikzpicture}[
%     font=\scriptsize,
%     >={Stealth[length=1.5mm]},
%     box/.style={draw, rounded corners=1.3pt, align=center, minimum height=6mm,
%                 inner sep=1.6pt, font=\scriptsize},
%     emb/.style={box, draw=methodArchInputLine, fill=methodArchInputBlock,
%                 minimum width=15mm},
%     norm/.style={box, draw=methodArchFrame, fill=methodArchPairBlock,
%                  minimum width=7mm},
%     sca/.style={box, draw=methodArchTimeLine, fill=methodArchTimeBlock,
%                 minimum width=8mm},
%     attn/.style={box, draw=methodArchFrame, fill=methodArchAttnBlock,
%                  minimum width=13mm},
%     mlp/.style={box, draw=methodArchFrame, fill=methodArchMlpBlock,
%                 minimum width=12mm},
%     tri/.style={box, draw=methodArchPairLine, fill=methodArchPairBlock,
%                 minimum width=13mm},
%     head/.style={box, draw=methodArchFrame, fill=methodArchAttnBlock,
%                  minimum width=14mm},
%     time/.style={box, draw=methodArchTimeLine, fill=methodArchTimeBlock,
%                  minimum width=15mm},
%     sum/.style={draw, circle, inner sep=0pt, minimum size=3.2mm, font=\tiny},
%     flow/.style={->, draw=methodArchArrow, line width=0.45pt},
%     wire/.style={draw=methodArchArrow, line width=0.45pt},
%     region/.style={rounded corners=4pt, draw, inner sep=5pt},
%     title/.style={font=\bfseries\footnotesize},
%   ]
%     % Condition-only input representations
%     \node[emb] (singlecond) {SingleCond};
%     \node[emb, below=13mm of singlecond] (paircond) {PairCond};
%     \node[left=4mm of singlecond, font=\small] (condition) {$\mathcal{C}$};
%     \draw[flow] (condition) -- (singlecond.west);
%     \draw[flow] (condition.south) |- (paircond.west);
%     \coordinate (singlebranch) at ([xshift=2.2mm]singlecond.east);
%     \node[circle, fill=methodArchArrow, inner sep=0pt, minimum size=1.2pt]
%       at (singlebranch) {};
%     \draw[flow] (singlebranch) |- ([yshift=3mm]paircond.north) -- (paircond.north);
%     \begin{scope}[on background layer]
%       \node[region, draw=methodArchInputLine, fill=methodArchInputBg,
%             fit=(singlecond)(paircond)] (conditionbg) {};
%     \end{scope}
%     \node[title, above=0.3mm of conditionbg.north] {Condition Embedder};

%     % Pair-only Pairmixer
%     \node[tri, right=6mm of paircond] (triout) {TriMulOut};
%     \node[sum, right=1.8mm of triout] (pairaddone) {$+$};
%     \node[tri, right=1.8mm of pairaddone] (triin) {TriMulIn};
%     \node[sum, right=1.8mm of triin] (pairaddtwo) {$+$};
%     \node[mlp, right=1.8mm of pairaddtwo] (transition) {Transition};
%     \node[sum, right=1.8mm of transition] (pairaddthree) {$+$};
%     \draw[flow] (paircond.east) -- (triout.west)
%       node[midway, above, font=\tiny] {$\mathbf{P}$};
%     \draw[flow] (triout) -- (pairaddone);
%     \draw[flow] (pairaddone) -- (triin);
%     \draw[flow] (triin) -- (pairaddtwo);
%     \draw[flow] (pairaddtwo) -- (transition);
%     \draw[flow] (transition) -- (pairaddthree);
%     \draw[flow] ([xshift=-2mm]triout.west) -- ++(0,-3.2mm)
%       -| (pairaddone.south);
%     \draw[flow] ([xshift=-1.6mm]triin.west) -- ++(0,-3.2mm)
%       -| (pairaddtwo.south);
%     \draw[flow] ([xshift=-1.6mm]transition.west) -- ++(0,-3.2mm)
%       -| (pairaddthree.south);
%     \coordinate (pairbottom) at ([yshift=-4.5mm]pairaddtwo.south);
%     \begin{scope}[on background layer]
%       \node[region, draw=methodArchPairLine, fill=methodArchPairBg,
%             fit=(triout)(pairaddthree)(pairbottom)] (pairbg) {};
%     \end{scope}
%     \node[title, above=0.3mm of pairbg.north] {Pairmixer ($\times 4$)};

%     % Noisy state enters only after Pairmixer
%     \node[sum] (singleadd) at
%       ([xshift=25mm]pairaddthree.east |- singlecond) {$+$};
%     \node[emb, above=3mm of singleadd] (coordembed) {Coord};
%     \node[emb, below=3mm of singleadd] (latticeembed) {Lattice};
%     \node[left=3mm of coordembed, font=\small] (coordinput) {$\mathbf{X}_t$};
%     \node[left=3mm of latticeembed, font=\small] (latticeinput)
%       {$\boldsymbol{\ell}_t$};
%     \draw[flow] (coordinput) -- (coordembed);
%     \draw[flow] (latticeinput) -- (latticeembed);
%     \draw[flow] (coordembed.south) -- (singleadd.north);
%     \draw[flow] (latticeembed.north) -- (singleadd.south);
%     \draw[flow] (singlecond.east) -- (singleadd.west)
%       node[pos=0.55, above, font=\tiny] {$\mathbf{S}$};
%     \begin{scope}[on background layer]
%       \node[region, draw=methodArchInputLine, fill=methodArchInputBg,
%             fit=(coordembed)(latticeembed)(singleadd)] (noisybg) {};
%     \end{scope}
%     \node[title, above=0.3mm of noisybg.north] {Noisy-State Embedder};

%     % DiT trunk
%     \coordinate (trunkstart) at ([xshift=13mm]singleadd.east);
%     \node[norm, anchor=west] (normone) at (trunkstart) {LN};
%     \node[sca, right=1.6mm of normone] (shiftone) {Scale\\Shift};
%     \node[attn, right=1.6mm of shiftone] (attention) {Attention};
%     \node[sca, right=1.6mm of attention] (gateone) {Scale};
%     \node[sum, right=1.6mm of gateone] (singleaddone) {$+$};
%     \node[norm, right=2.4mm of singleaddone] (normtwo) {LN};
%     \node[sca, right=1.6mm of normtwo] (shifttwo) {Scale\\Shift};
%     \node[mlp, right=1.6mm of shifttwo] (feedforward) {MLP};
%     \node[sca, right=1.6mm of feedforward] (gatetwo) {Scale};
%     \node[sum, right=1.6mm of gatetwo] (singleaddtwo) {$+$};
%     \draw[flow] (singleadd) -- (normone.west)
%       node[pos=0.52, above, font=\tiny] {$\widetilde{\mathbf{S}}$};
%     \draw[flow] (normone) -- (shiftone);
%     \draw[flow] (shiftone) -- (attention);
%     \draw[flow] (attention) -- (gateone);
%     \draw[flow] (gateone) -- (singleaddone);
%     \draw[flow] (singleaddone) -- (normtwo);
%     \draw[flow] (normtwo) -- (shifttwo);
%     \draw[flow] (shifttwo) -- (feedforward);
%     \draw[flow] (feedforward) -- (gatetwo);
%     \draw[flow] (gatetwo) -- (singleaddtwo);
%     \coordinate (skipone) at ([xshift=-1mm]normone.west);
%     \coordinate (skiptwo) at ($(singleaddone.east)!0.45!(normtwo.west)$);
%     \draw[flow] (skipone) -- ++(0,5mm) -| (singleaddone.north);
%     \draw[flow] (skiptwo) -- ++(0,6.3mm) -| (singleaddtwo.north);
%     \draw[flow] (pairaddthree.east) -| (attention.south)
%       node[pos=0.28, below, font=\tiny] {pair bias $\mathbf{P}^{\star}$};

%     \node[sca, minimum width=16mm, below=9mm of singleaddone]
%       (modulation) {Modulation};
%     \coordinate (modulationbus) at ($(modulation.north)+(0,4mm)$);
%     \draw[wire] (modulation.north) -- (modulationbus);
%     \foreach \target in {shiftone,gateone,shifttwo,gatetwo} {
%       \draw[flow] (modulationbus) -| (\target.south);
%     }
%     \begin{scope}[on background layer]
%       \node[region, draw=methodArchFrame, fill=methodArchTrunkBg,
%             fit=(normone)(singleaddtwo)(modulation)] (trunkbg) {};
%     \end{scope}
%     \node[title, above=0.3mm of trunkbg.north] {DiT Trunk ($\times 16$)};

%     % Time modulation
%     \node[time] (timeembed) at ([yshift=-13mm]noisybg.south)
%       {SinuEmb $\to$ MLP};
%     \node[left=4mm of timeembed, font=\small] (timeinput) {$t$};
%     \draw[flow] (timeinput) -- (timeembed);
%     \draw[flow] (timeembed.east) -| (modulation.south)
%       node[pos=0.22, above, font=\tiny] {$c_t$};
%     \begin{scope}[on background layer]
%       \node[region, draw=methodArchTimeLine, fill=methodArchTimeBg,
%             fit=(timeembed)] (timebg) {};
%     \end{scope}
%     \node[title, above=0.3mm of timebg.north] {Time Embedder};

%     % Coordinate and lattice endpoint heads
%     \node[head] (coordhead) at ($(singleaddtwo.east)+(15mm,4mm)$)
%       {CoordHead};
%     \node[head, below=2.5mm of coordhead] (latticehead) {LatticeHead};
%     \draw[flow] (singleaddtwo) -- ($(singleaddtwo.east)+(7mm,0)$)
%       coordinate (headbranch);
%     \draw[flow] (headbranch) |- (coordhead.west);
%     \draw[flow] (headbranch) |- (latticehead.west);
%     \node[right=4mm of coordhead, font=\small] (coordoutput)
%       {$\widehat{\mathbf{X}}_1$};
%     \node[right=4mm of latticehead, font=\small] (latticeoutput)
%       {$\widehat{\boldsymbol{\ell}}_1$};
%     \draw[flow] (coordhead) -- (coordoutput);
%     \draw[flow] (latticehead) -- (latticeoutput);
%     \begin{scope}[on background layer]
%       \node[region, draw=methodArchFrame, fill=methodArchHeadBg,
%             fit=(coordhead)(latticehead)] (headbg) {};
%     \end{scope}
%     \node[title, above=0.3mm of headbg.north] {Endpoint Heads};
%   \end{tikzpicture}%
%   }
%   \caption{\textbf{Final crystal structure generator architecture.}
%   The condition initializes single and pair representations $\mathbf{S}$ and
%   $\mathbf{P}$. Four Pairmixer blocks update only $\mathbf{P}$ using outgoing and
%   incoming triangle multiplication followed by a transition; $\mathbf{S}$ bypasses
%   Pairmixer unchanged. Embeddings of the noisy coordinates $\mathbf{X}_t$ and
%   lattice latent $\boldsymbol{\ell}_t$ are added to the single track only after
%   Pairmixer. The refined condition-only pair track $\mathbf{P}^{\star}$ biases each
%   DiT attention layer, while the time embedding modulates the trunk. Separate heads
%   predict $\widehat{\mathbf{X}}_1$ and $\widehat{\boldsymbol{\ell}}_1$. The final
%   configuration contains neither a geometry embedding module nor single-track
%   Pairmixer updates.}
%   \label{fig:method-architecture}
% \end{figure}

% Final crystal structure generator architecture.
% Figure-local palette matches the architecture-ablation diagram.
\definecolor{methodArchInputBg}{HTML}{E9F0F5}
\definecolor{methodArchInputBlock}{HTML}{AFC9DA}
\definecolor{methodArchInputLine}{HTML}{5D8198}
\definecolor{methodArchTimeBg}{HTML}{FCF5E6}
\definecolor{methodArchTimeBlock}{HTML}{F6DFB1}
\definecolor{methodArchTimeLine}{HTML}{B7892D}
\definecolor{methodArchPairBg}{HTML}{E3EEE9}
\definecolor{methodArchPairBlock}{HTML}{C5DED1}
\definecolor{methodArchPairLine}{HTML}{5F8D7A}
\definecolor{methodArchTrunkBg}{HTML}{E7DCE3}
\definecolor{methodArchHeadBg}{HTML}{F8EBE9}
\definecolor{methodArchAttnBlock}{HTML}{EAC8B1}
\definecolor{methodArchMlpBlock}{HTML}{CDB6C5}
\definecolor{methodArchArrow}{HTML}{2D2D2D}
\definecolor{methodArchFrame}{HTML}{30343B}

\begin{figure}[t]
  \centering
  \resizebox{\textwidth}{!}{%
  \begin{tikzpicture}[
    font=\scriptsize,
    >={Stealth[length=1.5mm]},
    box/.style={draw, rounded corners=1.3pt, align=center, minimum height=6mm,
                inner sep=1.6pt, font=\scriptsize},
    emb/.style={box, draw=methodArchInputLine, fill=methodArchInputBlock,
                minimum width=15mm},
    norm/.style={box, draw=methodArchFrame, fill=methodArchPairBlock,
                 minimum width=7mm},
    sca/.style={box, draw=methodArchTimeLine, fill=methodArchTimeBlock,
                minimum width=8mm},
    attn/.style={box, draw=methodArchFrame, fill=methodArchAttnBlock,
                 minimum width=13mm},
    mlp/.style={box, draw=methodArchFrame, fill=methodArchMlpBlock,
                minimum width=12mm},
    tri/.style={box, draw=methodArchPairLine, fill=methodArchPairBlock,
                minimum width=13mm},
    head/.style={box, draw=methodArchFrame, fill=methodArchAttnBlock,
                 minimum width=14mm},
    time/.style={box, draw=methodArchTimeLine, fill=methodArchTimeBlock,
                 minimum width=15mm},
    sum/.style={draw, circle, inner sep=0pt, minimum size=3.2mm,
                font=\scriptsize},
    flow/.style={->, draw=methodArchArrow, line width=0.45pt},
    wire/.style={draw=methodArchArrow, line width=0.45pt},
    bdot/.style={circle, fill=methodArchArrow, inner sep=0pt,
                 minimum size=2.2pt},
    region/.style={rounded corners=4pt, draw, inner sep=5pt},
    title/.style={font=\bfseries\footnotesize},
  ]
    % Condition-only input representations
    \node[emb] (singlecond) {SingleCond};
    \node[emb, below=17mm of singlecond] (paircond) {PairCond};
    \coordinate (inputcolumn) at ([xshift=-17mm]singlecond);
    \node[font=\small, text=methodArchFrame, inner sep=1pt]
      (condition) at (inputcolumn |- singlecond) {$\mathcal{C}$};
    \coordinate (conditionbranch) at ([xshift=-3mm]singlecond.west);
    \node[bdot] at (conditionbranch) {};
    \draw[wire] (condition) -- (conditionbranch);
    \draw[flow] (conditionbranch) -- (singlecond.west);
    \draw[flow] (conditionbranch) |- (paircond.west);
    \coordinate (singlebranch) at ([xshift=3mm]singlecond.east);
    \coordinate (pairtop) at ([yshift=3.5mm]paircond.north);
    \coordinate (pairbend) at (singlebranch |- pairtop);
    \node[bdot] at (singlebranch) {};
    \draw[flow] (singlebranch) -- (pairbend) -- (pairtop) -- (paircond.north);
    \coordinate (inputtop) at ([yshift=2mm]singlecond.north);
    \begin{scope}[on background layer]
      \node[region, draw=methodArchInputLine, fill=methodArchInputBg,
            fit=(inputtop)(conditionbranch)(singlecond)(paircond)(singlebranch)
                (pairbend)(pairtop)] (conditionbg) {};
    \end{scope}
    \node[title, above=0.7mm of conditionbg.north] (conditiontitle)
      {Input Embedder};

    % Pair-only Pairmixer
    \node[tri, right=20mm of paircond] (triout) {TriMulOut};
    \node[sum, right=2.6mm of triout] (pairaddone) {$+$};
    \node[tri, right=3mm of pairaddone] (triin) {TriMulIn};
    \node[sum, right=2.6mm of triin] (pairaddtwo) {$+$};
    \node[mlp, right=3mm of pairaddtwo] (transition) {Transition};
    \node[sum, right=2.6mm of transition] (pairaddthree) {$+$};
    \draw[flow] (paircond.east) -- (triout.west);
    \draw[flow] (triout) -- (pairaddone);
    \draw[flow] (pairaddone) -- (triin);
    \draw[flow] (triin) -- (pairaddtwo);
    \draw[flow] (pairaddtwo) -- (transition);
    \draw[flow] (transition) -- (pairaddthree);
    \coordinate (pairresone) at ([xshift=-4mm]triout.west);
    \coordinate (pairrestwo) at ([xshift=-1.5mm]triin.west);
    \coordinate (pairresthree) at ([xshift=-1.5mm]transition.west);
    \foreach \point in {pairresone,pairrestwo,pairresthree} {
      \node[bdot] at (\point) {};
    }
    \draw[flow] (pairresone) -- ++(0,-5.5mm) -| (pairaddone.south);
    \draw[flow] (pairrestwo) -- ++(0,-5.5mm) -| (pairaddtwo.south);
    \draw[flow] (pairresthree) -- ++(0,-5.5mm) -| (pairaddthree.south);
    \coordinate (pairtopbound) at ([yshift=2mm]triout.north);
    \coordinate (pairbottom) at ([yshift=-7mm]triout.south);
    \begin{scope}[on background layer]
      \node[region, draw=methodArchPairLine, fill=methodArchPairBg,
            fit=(pairresone)(triout)(pairaddthree)(pairtopbound)(pairbottom)]
            (pairbg) {};
    \end{scope}
    \node[title, above=0.7mm of pairbg.north] (pairtitle)
      {Pairmixer ($\times 4$)};

    % Noisy state enters only after Pairmixer
    \node[sum] (singleadd) at
      ([xshift=24mm]pairaddthree.east |- singlecond) {$+$};
    \node[emb, above=3mm of singleadd] (coordembed) {Coord};
    \node[emb, below=3mm of singleadd] (latticeembed) {Lattice};
    \node[left=4mm of coordembed, font=\small, text=methodArchFrame,
          inner sep=1pt] (coordinput) {$\mathbf{X}_t$};
    \node[left=4mm of latticeembed, font=\small, text=methodArchFrame,
          inner sep=1pt] (latticeinput) {$\boldsymbol{\ell}_t$};
    \draw[flow] (coordinput) -- (coordembed);
    \draw[flow] (latticeinput) -- (latticeembed);
    \draw[flow] (coordembed.south) -- (singleadd.north);
    \draw[flow] (latticeembed.north) -- (singleadd.south);
    \draw[flow] (singlecond.east) -- (singleadd.west);
    \begin{scope}[on background layer]
      \node[region, draw=methodArchInputLine, fill=methodArchInputBg,
            fit=(coordembed)(latticeembed)(singleadd)] (noisybg) {};
    \end{scope}
    \node[title, above=0.7mm of noisybg.north] (noisytitle)
      {Noisy Input Embedder};

    % Time-conditioned DiT trunk
    \coordinate (trunkstart) at ([xshift=16mm]coordembed.east |- singlecond);
    \node[norm, anchor=west] (normone) at (trunkstart) {LN};
    \node[sca, right=1.6mm of normone] (shiftone) {Scale\\Shift};
    \node[attn, right=1.6mm of shiftone] (attention) {Attention};
    \node[sca, right=1.6mm of attention] (gateone) {Scale};
    \node[sum, right=1.6mm of gateone] (singleaddone) {$+$};
    \node[norm, right=3mm of singleaddone] (normtwo) {LN};
    \node[sca, right=1.6mm of normtwo] (shifttwo) {Scale\\Shift};
    \node[mlp, right=1.6mm of shifttwo] (feedforward) {MLP};
    \node[sca, right=1.6mm of feedforward] (gatetwo) {Scale};
    \node[sum, right=1.6mm of gatetwo] (singleaddtwo) {$+$};
    \draw[flow] (singleadd) -- (normone.west);
    \draw[flow] (normone) -- (shiftone);
    \draw[flow] (shiftone) -- (attention);
    \draw[flow] (attention) -- (gateone);
    \draw[flow] (gateone) -- (singleaddone);
    \draw[flow] (singleaddone) -- (normtwo);
    \draw[flow] (normtwo) -- (shifttwo);
    \draw[flow] (shifttwo) -- (feedforward);
    \draw[flow] (feedforward) -- (gatetwo);
    \draw[flow] (gatetwo) -- (singleaddtwo);
    \coordinate (skipone) at ([xshift=-4.5mm]normone.west);
    \coordinate (skiptwo) at ([xshift=-1.5mm]normtwo.west);
    \node[bdot] at (skipone) {};
    \node[bdot] at (skiptwo) {};
    \draw[flow] (skipone) -- ++(0,5.5mm) -| (singleaddone.north);
    \draw[flow] (skiptwo) -- ++(0,5.5mm) -| (singleaddtwo.north);
    \draw[flow] (pairaddthree.east) -| (attention.south);

    \node[sca, minimum width=16mm, below=8mm of singleaddone]
      (modulation) {Modulation};
    \coordinate (modulationbus) at ($(modulation.north)+(0,2.5mm)$);
    \draw[wire] (modulation.north) -- (modulationbus);
    \node[bdot] at (modulationbus) {};
    \foreach \target in {shiftone,gateone,shifttwo,gatetwo} {
      \draw[flow] (modulationbus) -| (\target.south);
    }
    \coordinate (trunktop) at ([yshift=7mm]normone.north);
    \coordinate (trunkleft) at ([xshift=-1.5mm]skipone);
    \coordinate (trunkbottom) at ([yshift=-2.5mm]modulation.south);
    \begin{scope}[on background layer]
      \node[region, draw=methodArchFrame, fill=methodArchTrunkBg,
            fit=(normone)(singleaddtwo)(modulation)(trunktop)(trunkleft)
                (trunkbottom)] (trunkbg) {};
    \end{scope}
    \node[title, above=0.7mm of trunkbg.north] (trunktitle)
      {Transformer Trunk ($\times 16$)};

    % Time Embedder below the Input Embedder, as in Figure 8(b)
    \node[time] (timeembed) at ([yshift=-13.5mm]conditionbg.south)
      {SinuEmb $\to$ MLP};
    \node[font=\small, text=methodArchFrame, inner sep=1pt]
      (timeinput) at (inputcolumn |- timeembed) {$t$};
    \draw[flow] (timeinput) -- (timeembed);
    \draw[flow] (timeembed.east) -| (modulation.south);
    \begin{scope}[on background layer]
      \node[region, draw=methodArchTimeLine, fill=methodArchTimeBg,
            fit=(timeembed)] (timebg) {};
    \end{scope}
    \node[title, above=0.7mm of timebg.north] (timetitle) {Time Embedder};

    % Coordinate and lattice endpoint heads
    \node[head] (coordhead) at ($(singleaddtwo.east)+(18mm,5mm)$)
      {CoordHead};
    \node[head, below=3.4mm of coordhead] (latticehead) {LatticeHead};
    \draw[wire] (singleaddtwo) -- ($(singleaddtwo.east)+(7mm,0)$)
      coordinate (headbranch);
    \node[bdot] at (headbranch) {};
    \draw[flow] (headbranch) |- (coordhead.west);
    \draw[flow] (headbranch) |- (latticehead.west);
    \node[right=4mm of coordhead, font=\small] (coordoutput)
      {$\widehat{\mathbf{X}}_1$};
    \node[right=4mm of latticehead, font=\small] (latticeoutput)
      {$\widehat{\boldsymbol{\ell}}_1$};
    \draw[flow] (coordhead) -- (coordoutput);
    \draw[flow] (latticehead) -- (latticeoutput);
    \begin{scope}[on background layer]
      \node[region, draw=methodArchFrame, fill=methodArchHeadBg,
            fit=(headbranch)(coordhead)(latticehead)] (headbg) {};
    \end{scope}
    \node[title, above=0.7mm of headbg.north] (headtitle) {Crystal Heads};

    % Complete denoiser frame
    \begin{scope}[on background layer]
      \node[region, draw=methodArchFrame, fill=none, inner sep=7pt,
            fit=(conditionbg)(conditiontitle)(pairbg)(pairtitle)
                (noisybg)(noisytitle)(trunkbg)(trunktitle)
                (headbg)(headtitle)(timebg)(timetitle)
                (condition)(timeinput)(coordoutput)] (denoiserbg) {};
    \end{scope}
  \end{tikzpicture}%
  }
    \caption{\textbf{Model architecture of \model{}.}
    The condition $\mathcal{C}$ is embedded into single and pair representations,
    with the pair representation refined by four \textsc{Pairmixer}
    blocks~\citep{ouyangzhang2025triangle}. Because this computation depends only on
    $\mathcal{C}$, it can be cached across denoising steps. Noisy coordinate
    $\mathbf{X}_t$ and lattice $\boldsymbol{\ell}_t$ embeddings are then added to the
    single representation and processed by a DiT trunk, with the refined pair
    representation providing attention biases throughout. Separate heads predict the
    coordinate and lattice endpoints $\widehat{\mathbf{X}}_1$ and
    $\widehat{\boldsymbol{\ell}}_1$.}
    \label{fig:method-architecture}
\end{figure}


\textbf{Condition embedding.}
The input embedder embeds the condition
$\mathcal{C}=(\mathcal{G},\mathbf{r},Z,\mathcal{O})$ into a per-atom single
representation $\mathbf{S}_{\mathcal C}$ and an initial per-atom-pair
representation $\mathbf{P}^{(0)}$:
\begin{equation*}
\mathbf{S}_{\mathcal C}=E_{\mathrm{single}}(\mathcal C),
\qquad
\mathbf{P}^{(0)}=E_{\mathrm{pair}}(\mathbf{S}_{\mathcal C},\mathcal C),
\end{equation*}
where $E_{\mathrm{single}}$ and $E_{\mathrm{pair}}$ are learned single and pair
embedding networks. The single representation combines atom identity,
periodic-table descriptors, formal charge, optional template coordinates,
chirality, and space group. The pair representation combines projected single
features with intramolecular bond types, bond stereochemistry, template
displacements, and template distances. Availability masks and learned null states
distinguish unavailable optional conditions from provided values.

\textbf{Pairmixer.}
Pairmixer is a streamlined alternative to Pairformer that retains incoming and
outgoing triangle multiplication and pair transitions while omitting triangle
attention. Each block updates the current pair representation $\mathbf{P}$ through
sequential residual operations,
\begin{equation*}
\mathbf{P}\leftarrow\mathbf{P}
+\operatorname{TriMul}_{\mathrm{out}}(\mathbf{P}),\qquad
\mathbf{P}\leftarrow\mathbf{P}
+\operatorname{TriMul}_{\mathrm{in}}(\mathbf{P}),\qquad
\mathbf{P}\leftarrow\mathbf{P}
+\operatorname{Transition}(\mathbf{P}),
\end{equation*}
where $\operatorname{TriMul}_{\mathrm{out}}$ and
$\operatorname{TriMul}_{\mathrm{in}}$ denote outgoing and incoming triangle
multiplication, respectively. The transition is a gated feed-forward network
applied independently to each atom pair. Pairmixer leaves the single representation unchanged and produces a refined pair representation $\mathbf{P}^{\star}$.

\textbf{Noisy-state injection.}
The noisy crystal state is injected into the single representation
$\mathbf{S}_{\mathcal C}$. Coordinate features are encoded using Fourier features~\citep{tancik2020fourier}, while the lattice is embedded
globally and broadcast to all atoms. The resulting single representation is
\begin{equation*}
\mathbf{S}_t
=\mathbf{S}_{\mathcal C}
+E_X(\mathbf{X}_t)
+\operatorname{Broadcast}\!\left(E_L(\boldsymbol{\ell}_t)\right),
\end{equation*}
where $E_X$ and $E_L$ denote the coordinate and lattice embedding networks,
respectively. A sinusoidal embedding of the flow time $t$ provides a separate
global condition to the transformer.

\textbf{DiT trunk.}
The trunk follows the DiT architecture, alternating between self-attention and
feed-forward updates to the single representation. In each attention layer, the
refined pair representation $\mathbf{P}^{\star}$ is projected to a
head-specific bias and added to the attention logits:
\begin{equation*}
a_{ij}^{(h)}
=
\frac{\left\langle\mathbf{q}_i^{(h)},\mathbf{k}_j^{(h)}\right\rangle}
{\sqrt{d_h}}
+b_h\!\left(\mathbf{P}_{ij}^{\star}\right),
\end{equation*}
where $\mathbf{q}_i^{(h)}$ and $\mathbf{k}_j^{(h)}$ are the query and key
vectors for atoms $i$ and $j$ in attention head $h$, $d_h$ is the head
dimension, and $b_h$ maps the pair representation to a scalar attention bias.
The time embedding modulates the normalization and residual gates of both
self-attention and feed-forward updates~\citep{peebles2023scalable}.

\textbf{Prediction heads.}
The DiT output $\mathbf{S}^{\mathrm{out}}$ is decoded by separate coordinate and
lattice endpoint heads. For the set $\mathcal V$ of valid atoms, the coordinate
head predicts an atomwise vector $\mathbf r_i$ and subtracts its mean over valid
atoms, while the lattice head applies an MLP to the mean-pooled single
representation:
\begin{equation*}
\mathbf{r}_i
= W_X \operatorname{Norm}(\mathbf{S}^{\mathrm{out}}_i),\qquad
\widehat{\mathbf{X}}_{1,i}
= \mathbf{r}_i
- \frac{1}{|\mathcal V|}\sum_{j\in\mathcal V}\mathbf{r}_j,\qquad
\widehat{\boldsymbol{\ell}}_1
= \operatorname{MLP}\!\left(
\frac{1}{|\mathcal V|}\sum_{i\in\mathcal V}
\operatorname{Norm}(\mathbf{S}^{\mathrm{out}}_i)
\right),
\end{equation*}
where $W_X$ is the coordinate projection and
$\widehat{\mathbf{X}}_1$ and $\widehat{\boldsymbol{\ell}}_1$ denote the predicted
coordinate and lattice endpoints.